\documentclass[11pt]{article}
\usepackage{amsmath,amssymb,amsthm,geometry,hyperref}
\geometry{margin=1in}
\newtheorem{theorem}{Theorem}
\newtheorem{lemma}{Lemma}
\newtheorem{proposition}{Proposition}
\newcommand{\Z}{\mathbb Z}
\newcommand{\Q}{\mathbb Q}
\newcommand{\C}{\mathbb C}
\newcommand{\pr}{\operatorname{pr}}
\newcommand{\im}{\operatorname{im}}
\newcommand{\Tr}{\operatorname{tr}}
\newcommand{\calH}{\mathcal H}
\newcommand{\calD}{\mathcal D}
\title{Integral gluing and ordered fillings for the actual
four-variable prepared germs in manuscript family 059}
\author{Campaign 9000059: Engineer proof submitted for independent review}
\date{November 5, 2026}
\begin{document}
\maketitle

\begin{abstract}
We identify the integral local Milnor lattice and its ordered filling
form for each of the two actual prepared germs in M45. The comparison
with a local page is an integral relative-homology isomorphism, not a
rank argument. A root-disk calculation, natural under every braid
and hence under the actual graph covering, determines the mod-four
extension. We prove both cross-term vanishings and the relative signs
by geometric fillings. A direct quantitative isolation argument
avoids curve selection in the degenerating parameter. Compact exact
certificates check the root-disk index, every four-sheet permutation
over \(\Z/4\), the cap-boundary signs, and the normal-intersection
determinants. These computations support, but do not replace, the
geometric proof. This is not a Lean/kernel verification. Independent
review and novelty assessment remain outstanding; neither manuscript's
ambient-equivalence theorem is decided here.
\end{abstract}

\section{Question, inputs, and conclusion}

M45 denotes the staged September 27, 2026 manuscript
\emph{Ambiently homeomorphic isolated hypersurface germs in
\(\C^4\) with multiplicities four and five}.
We address its \texttt{thm:tube}, not its subsequent arithmetic or
classification assertions. The two graph maps are
\[
\begin{split}
h(X,T)&=T(X-T^2)(X-2T^2),\\
q_1(u,w)&=(u-w-w^8)(u-w-2w^8)(u-w^2),\qquad q_2=q_1-u^2,\\
\pi_0(u,w)&=(u^2+q_1(u,w),u),\\
\pi_1(x,w)&=(x,u(u-w^2)q_2(u,w)),\qquad u=x+w .
\end{split}
\]
For either graph let \(d=h\circ\pi\), and use the exact monic
Weierstrass polynomial \(R(p,u)=u^4+\sum_{j<4}A_j(p)u^j\)
of its graph equation. No analytic unit is removed from the
definition of the fourfold or its function. Set
\[
 R_s(p,w)=s^4R(p,w/s),\quad C_s=\{yz=R_s(p,w)\},\quad
 f_s=h(p)+w(y^a+z^b).
\]
Here \(s>0\), \(a,b\) are the manuscript's common fixed integers, and
\[
 \alpha=\frac{4b}{a+b},\quad\beta=\frac{4a}{a+b},\quad
 \epsilon_a=\frac{(a-2)\alpha}{4}-4>0,\quad
 \epsilon_b=\frac{(b-2)\beta}{4}-4>0.                 \tag{1}
\]
In particular \(a,b>2\). The common arithmetic family in Section 03
is contained in this range once its integer parameter is sufficiently
large. Nothing below requires a further arithmetic hypothesis.

The previously independently reviewed claims \texttt{C059BASE},
\texttt{C059GRADIENT}, and \texttt{C059APPLICATION} give:
\[
\begin{gathered}
\deg\pi=4,\quad A_j(0)=0,\quad
R=0\text{ is smooth with a unit base derivative},\\
\mu(h)=6,\quad\mu(d)=51,\quad
\mu(f_s)=24ab+6a+6b+51.
\end{gathered}                                      \tag{2}
\]
The last equality and local isolation hold for each fixed \(s\ne0\).
They do not give uniform isolation on a fixed outer tube.
The reducedness and isolation of the two plane germs follow from
their finite Jacobian lengths. We reuse these reviewed inputs,
not the manuscript's plane-form isometry or its monodromy-sector
claims.

We use the ordinary local Milnor fibration and contractibility of its
closed disk total, in precisely \cite[Theorem 1.8 and Remark 1.9]{Massey}.
The isolated-singularity homotopy and Jacobian-number statements are
recorded in \cite[Section 1, p.~6, items (2)--(5)]{Massey}.
The actual versioned primary text was obtained and read, and is retained
beside this report. These statements concern ordinary local tubes:
their applicability to our custom cuts is proved below.
All further gluing, braid, and sign arguments are supplied here.
In particular, no Sakamoto or Kauffman--Neumann product theorem is
assumed. An earlier landing-page inspection of Sakamoto was not
verification of a theorem.

Let \(\calH,\calD\) be compact plane pages; put
\(H=H_1(\calH;\Z)\), \(D=H_1(\calD;\Z)\).
The graph induces \(P_*=\pi_*:D\to H\) and full-preimage transfer
\(T_*=\pi^!:H\to D\). A loop may be moved off the finitely many
branch values. Its full inverse image defines the transfer.
Moving a loop across a branch value does not change this homology
class: pull back the generic oriented surface implementing the
homology, allowing its isolated branch points. Its boundary is
the difference of the full inverse-image cycles. Thus the transfer
is well defined, and sheet counting gives
\[
 P_*T_*=4I_H,\qquad K=\ker P_*,\qquad
 \pr(l)=l-\tfrac14T_*P_*l.                          \tag{3}
\]
Both plane pages are connected surfaces with boundary and have
free first homology. These facts also follow from the isolated
Milnor-fiber statement just cited.

On the cone \(V=\{yz=w^4\}\) put \(P=w(y^a+z^b)\).
Let \(Q\) be a compact small page of \(P\), \(L\) its weighted link,
\(E=H_1(Q;\Z)\), and let \(\tau:E\to H_1(L)=\Z/4\) be radial
projection. Orient the generator of \(H_1(L)\) by a positive
\(y\)-angle at a fixed nonzero \(w\).
The cover
\[
 (U,V)\longmapsto (U^4,V^4,UV),\qquad
 \widetilde P=UV(U^{4a}+V^{4b})
\]
is unramified off the vertex. Write
\(\Tr:E\to H_1(\widetilde Q;\Z)\) for its full lift and define
\[
 e(v,v')=\tfrac14\widetilde s(\Tr v,\Tr v').         \tag{4}
\]
The forms \(s_h,s_d,\widetilde s\) are ordered intersections of
fillings in their complex-oriented plane totals; the second
boundary is on the positive neighboring page.

\paragraph{Chain conventions.}
All homology below is integral unless coefficients are specified.
The elementary exact sequences used in the proof come from free
chain complexes, so they retain torsion. For a collared union
\(A\cup B\), subdivision gives the short exact chain sequence
\[
0\longrightarrow C_*(A\cap B)
\xrightarrow{c\mapsto(c,-c)}C_*(A)\oplus C_*(B)
\xrightarrow{(a,b)\mapsto a+b}C_*(A\cup B)\longrightarrow0;
\]
the relative version is obtained by quotienting the corresponding
page subcomplexes. Taking cycles and boundaries gives the
Mayer--Vietoris sequences used here, including (23) and (26).
Products use the boundary rule
\(\partial(A\times B)=\partial A\times B+
(-1)^{\dim A}A\times\partial B\).
A compact connected surface with nonempty boundary has a graph
spine. In products with either plane page the cellular complex
therefore has a free degree-zero part and a free degree-one part,
both with zero differential; the product homology formulas used
below follow by tensoring this complex. In particular no
unmentioned torsion product occurs.
The quotient cone link has the usual two-solid-torus description
obtained from \(|U|\le|V|\) and \(|V|\le|U|\) on its covering
three-sphere. Its meridian attachment wraps the positive \(y\)-circle
four times. Its one-cell-per-dimension complex consequently has
\(d_2=4\), \(d_1=d_3=0\); this also computes \(H_2(L)=0\)
without a duality assertion.
Finally the coefficient sequence is multiplication by four on
these free complexes; its connecting map is precisely
\([\overline C]\mapsto[\partial C/4]\).

\begin{theorem}[Actual integral tube identification]\label{thm:main}
For either actual prepared graph, every fixed pair satisfying (1),
and all sufficiently small \(s>0\), the origin is the only critical
point in a fixed small outer cut. The ordinary local Milnor page
has, in the coordinates constructed below, exactly the integral image
\[
 \Lambda_\pi=
 \left\{(\pr(l),z):l\in D,\ z\in H\otimes E,\
       P_*l\bmod4H=(1\otimes\tau)z\right\}
 \subset K_\Q\oplus(H\otimes E)_\Q .                \tag{5}
\]
Its ordered filling form is the restriction of
\[
              -s_d|_K\ \perp\ (s_h\otimes e).       \tag{6}
\]
There is no unmeasured finite index or torsion quotient in (5).
\end{theorem}

\section{Uniform isolation without a limiting-arc argument}

The following estimate makes the passage from a germ to a fixed
outer representative quantitative. Throughout, constants may depend
on the fixed graph and on \(a,b\), but not on \(s\).
With \(r=\max(|X|,|T|)\) in a sufficiently small base ball,
\[
             \max(|h_X|,|h_T|)\ge c r^4.           \tag{7}
\]
Here is an elementary proof, rather than an appeal to an arc criterion.
If \(|X|\ge20|T|^2\), then
\(|h_T|\ge(1-9/20-10/400)|X|^2\), which bounds \(r^4\)
below up to a constant. Otherwise put \(v=X/T^2\), \(|v|\le20\).
For \(A=2v-3\), \(B=v^2-9v+10\) the exact identity
\[
                  4B-(2v-15)A=-5                  \tag{8}
\]
gives \(\max(|A|,|B|)\ge5/59\). Since
\(h_X=T^3A,\ h_T=T^4B\), shrinking the ball proves (7).
The cases \(T=0\) are included in the first case.

Scale \(w=su,\ y=s^\alpha Y,\ z=s^\beta Z\), and put
\(\xi=s^{1+a\alpha}>0\). The constraint is \(YZ=R(p,u)\).
A critical point satisfies, for a complex multiplier \(\chi\),
\[
\begin{aligned}
 h_p&=-\chi R_p,&\xi(Y^a+Z^b)&=-\chi R_u,\\
 \xi a uY^{a-1}&=\chi Z,&
 \xi b uZ^{b-1}&=\chi Y.
\end{aligned}                                      \tag{9}
\]
When \(uYZ\ne0\), also \(\chi\ne0\). Multiplying the last two
equations by \(Y,Z\) and using the second gives
\[
 aY^a=bZ^b,\qquad
 ab\,uR_u+(a+b)R=0.                                \tag{10}
\]
The second polynomial has leading coefficient \(4ab+a+b\);
its coefficient of \(u^j\), \(j<4\), is
\((jab+a+b)A_j(p)=O(r)\).
Its roots satisfy \(|u|\le C r^{1/4}\).
Indeed a root with \(|u|\ge1\) would satisfy
\(|u|^4\le C r\sum_{j<4}|u|^j\le4Cr|u|^3\),
impossible for small \(r\); for \(|u|\le1\) the same equation
gives \(|u|^4\le4Cr\).
It follows that
\[
 |R|\le C r,\quad |Y|\le C r^{\alpha/4},\quad
 |Z|\le C r^{\beta/4},\quad |\chi|\ge c' r^4.        \tag{11}
\]
The last bound uses (7) and the boundedness of \(R_p\) on this
now bounded \((p,u)\)-set. Thus (9) implies simultaneously
\[
 \frac{|Z|}{|Y|}\le C_a\xi r^{1/4+\epsilon_a},\qquad
 \frac{|Y|}{|Z|}\le C_b\xi r^{1/4+\epsilon_b}.       \tag{12}
\]
On a fixed sufficiently small base ball, both right sides are
less than one for sufficiently small \(s\), a contradiction.
If \(p=0\), (10) instead gives \(u=0\), already a contradiction
to \(uYZ\ne0\).

If exactly one of \(Y,Z\) is zero, (9) forces \(\chi=0\),
and its \(u\)-equation forces the other to vanish.
If \(u=0\) and \(YZ\ne0\), then \(\chi=0\), \(h_p=0\), and
the constraint contradicts \(R(0,0)=0\).
Finally, when \(Y=Z=0\), the remaining equations say exactly that
\(h|_{R=0}=d\) is critical. Its only critical point in the selected
root representative is the origin by (2).
This proves uniform isolation, without assuming that \(w/s\)
was bounded in advance and without using local isolation to
assert a uniform statement.

\section{Representatives, continuations, and the local comparison}

We give the analytic details needed by the subsequent chain
calculation; no homology model is presumed to be a local page.

\subsection{Base cuts and boundary ranks}

Choose a sufficiently small round base ball \(\mathcal B\) and a
proper finite graph representative over it. The zero curves of \(h\)
and \(d\) are reduced and isolated. Their finitely many punctured
branches meet the base sphere, and its inverse image on the root
surface, transversely when the ball is sufficiently small. To see the
latter assertion directly, on a root-branch parametrization the
nonconstant analytic map \(p(t)\) has
\(|p(t)|^2=c|t|^{2m}(1+O(|t|))\), \(c>0\);
its radial derivative is positive. Along these branch annuli,
the respective functions are complex submersions in a normal
direction. Compactness supplies a small common value disk with
regular boundaries.

These finite graph cuts have the same acyclic totals as ordinary
plane tubes. Indeed, take an ordinary small ball inside a graph cut.
Between its boundary and the cut boundary, the zero curve is the
disjoint union of the preceding annuli. A field along each annulus
pointing from the inner boundary to the outer one extends in its
complex submersion neighborhood to a field preserving the function.
Outside these neighborhoods the compact intervening region has
the absolute value of the function bounded away from zero.
For a sufficiently small value disk the entire intervening tube
is therefore a value-preserving collar. Attaching it to the ordinary
contractible total changes no homology. This also supplies
boundary-respecting trivializations over compact regular-value regions.

Choose \(M>3\) with \(M^{a+b}>1+4\max(a,b)\), and put
\[
Y_t=\{|y|\le Mt^\alpha,\ |z|\le Mt^\beta,\ |w|\le t\}.
\]
Choose \(\rho\), then \(\varepsilon\), then
\(0<\delta<\varepsilon/4\), then \(\lambda\ll\rho\) so that
\(\sup_{Y_\lambda}|P|<\delta/4\), and finally \(s>0\)
so small that every root of \(R_s(p,\cdot)\) over the compact
base cut is in \(|w|<\lambda/4\).
Let
\[
\begin{split}
X_s&=\{p\in\mathcal B,\ |h(p)|\le\varepsilon,\
               (y,z,w)\in Y_\rho,\ yz=R_s(p,w)\},\\
\widehat{\mathcal T}_s&=X_s\cap\{|f_s|\le\delta\},\qquad
\mathcal T_s=X_s\cap\{f_s=c\},\quad |c|=\delta .
\end{split}                                                     \tag{13}
\]
The ambient \(C_s\) is smooth throughout these cuts: when \(y,z\)
both vanish, the root lies in the chosen graph neighborhood and
has a unit base derivative; elsewhere a vertical derivative is
nonzero.

We verify disk-wide rank two on every outer lateral stratum.
At \(s=0\), the \(y\)- and \(z\)-faces do not meet. On a
\(y\)-face, eliminate \(z=w^4/y\). Then
\[
 |z^b/y^a|\le M^{-(a+b)},\qquad
 P_w=y^a+(4b+1)z^b\ne0.                            \tag{14}
\]
The \(z\)-face is identical with the labels exchanged.
At a \(y/w\) corner,
\[
 |P|\ge \rho^{1+a\alpha}(M^a-M^{-b})>0,             \tag{15}
\]
and there is a corresponding \(z/w\) bound.
Take \(\varepsilon+\delta\) smaller than these fixed bounds.
On the \(w\)-face near \(P=0\),
\[
 dP(\partial_{\log y})=w(ay^a-bz^b)
                      =w(a+b)y^a\ne0.             \tag{16}
\]
Compactness gives a common small-value neighborhood.
In the nonvertex cone interior, the two derivatives of
\(\widetilde P\) cannot vanish together: if \(UV\ne0\) their
coefficient determinant is
\((4a+1)(4b+1)-1>0\); on either axis the transverse
derivative is nonzero.
The estimates on the fixed outer faces persist for small \(s\).
Their free vertical directions keep \(p\) fixed, so additional
base faces cause no loss of rank.

For completeness, possible degeneration towards the vertical
vertex on a stratum with no vertical face needs a separate argument.
Write
\(\varrho=\max(|y|^{1/\alpha},|z|^{1/\beta},|w|)\).
If \(\varrho/s\to\infty\), rescaling by \(\varrho\)
makes the vertical fiber converge smoothly to a nonvertex
point of the cone, and the normalized vertical derivative of \(P\)
is nonzero there. It is already a complex direction tangent
to every base cut. If \(\varrho/s\) is bounded, use
\((p,Y,Z,u)\) and (9). The functions converge in \(C^1\)
on compact sets to \(h\) on \(YZ=R(p,u)\).
Where \(Y\) or \(Z\) is nonzero, projection to \(p\)
is a submersion. Where both vanish, the root-graph boundary
submersion established above supplies the necessary rank.
The \(|h|=\varepsilon\) face is absent in this limit over
\(|f_s|\le\delta\), because \(P\to0\) and \(\delta<\varepsilon/4\).
The remaining base-sphere limits have \(p\ne0\).
This excludes every possible rank-loss sequence.
Thus (13) has proper disk-wide stratified submersive outer
boundary, including all its corners. Section 2 excludes interior
critical points other than the origin.

\subsection{Compatible inner and outer models}

Cut (13) at \(Y_\lambda\). Write \(W,M_{\rm ext},J\) for the
inner page, exterior page, and their bicollared interface; hats
denote disk totals. On the inner part replace \(h+P\) by
\(h+\theta P\), \(0\le\theta\le1\).
On its horizontal boundary
\[
                  3\delta/4\le |h|\le5\delta/4.    \tag{17}
\]
The preceding two-scale argument on this compact annulus, with
\(\theta\) an additional compact parameter, shows regularity
of these horizontal boundaries and their lateral intersections.
At \(s=0\), \(h\) is independent of the cone coordinate; for
bounded scaled variables the regular root restriction supplies
the same rank. Decreasing \(\lambda\) first controls the \(C^1\)
size of \(P\) away from the scaled root region, and decreasing
\(s\) controls the latter region.
Consequently the sublevel totals isotope, preserving the complex
value on the horizontal boundary. This does not claim that the
interior \(h\)-model is a submersion at value zero.

On the exterior, the ambient family and all boundary constraints
are smooth down to \(s=0\), so the analogous continuation gives
an isotopy of exterior total/page pairs. These assertions can be
proved by choosing local parameter lifts tangent to every active
face and solving the complex horizontal lift equation.
The conditions are affine in the lift. A partition of unity
patches the lifts; a cutoff extends them away from the horizontal
boundary. The compact domains make the flows complete.
This is the face-preserving submersion argument, not an inference
from homology ranks.

We specify an important relative normalization of these isotopies.
Let \(w_0=\lambda/2\), and choose an \(h\)-page trivialization
\(\Phi(p_0,q)\) near \(c\), respecting its base boundary.
For each \(p_0\in\calH_c\) and \(\phi\in\mathbb R/2\pi\Z\),
set
\[
 y=M\lambda^\alpha e^{i\phi},\quad w=w_0,\quad
 z=R_s(\Phi(p_0,q),w_0)/y,
\quad q+\theta P(y,z,w_0)=c.                        \tag{18}
\]
The last equation has a unique smooth solution near \(c\).
Its \(q\)-derivative differs from the identity by \(o(1)\):
every \(p\)-dependent coefficient of \(R_s(p,w_0)\)
contains a positive power of \(s\).
The root bound puts the circles strictly away from the other
vertical faces. The local parameter lifts solving (18) fix
\(p_0,w_0,\phi\); patch the general lifts to these near the
compact anchor trace. In the coordinates
\((p_0,q+\theta P,y,w,s,\theta)\), these anchored lifts
are parameter coordinate fields.
Thus the parameter rectangle has isotopic boundary
identifications relative to the anchored circles, with no
base-to-circle shear.
At \(s=0\) the full page interface is explicitly
\[
 (p_0,v)\longmapsto(\Phi(p_0,c-P(v)),v),\qquad v\in L.
                                                                    \tag{19}
\]
It is therefore \(\calH_c\times L\), with the prescribed
base label and positive \(y\)-circle.

\subsection{Acyclic total and integral local comparison}

In the inner \(h\)-model put
\[
 G=\calH_c\times\Delta_w,\qquad
 \calD=\{R_s=0\}\cap G .
\]
For a fixed \((p,w)\), balance the hyperbola by decreasing
\(|y|^2/\sigma^2-\sigma^2|z|^2\) linearly to zero while
keeping \(yz=R_s\) and the \(y\)-angle, where
\(\sigma\) is close to \(\lambda^{(\alpha-\beta)/2}\).
The final moduli are \(\sigma\sqrt{|R_s|}\) and
\(\sigma^{-1}\sqrt{|R_s|}\). The root bound and \(M>3\)
keep both strictly inside their cuts. Each modulus stays
between its initial and final values, also on the axes,
where the active modulus shrinks to zero.
This realizes a deformation retraction onto
\[
 X(G,\calD)=(G\times S^1)/
        (\{g\}\times S^1\text{ collapsed for }g\in\calD).
                                                                    \tag{20}
\]
The constant \(y\)-angle section is continuous across the roots.
It splits off \(G\); quotienting it gives
\((G/\calD)\wedge S^1\). Hence, integrally,
\[
 H_3(W)=K,\quad H_2(W)=H/P_*D,\quad
 H_3(W;\Z/4)\hookrightarrow D/4D.                  \tag{21}
\]
These follow from the relative long exact sequence of \((G,\calD)\);
there is no rank-only inference.

The same construction for disk totals has acyclic base and root
totals, so \(\widehat W\) is acyclic.
At \(s=0\), the weighted contraction
\[
 (X,T;y,z,w)\longmapsto
 (t^{2/5}X,t^{1/5}T;\
 t^{\alpha/d_P}y,t^{\beta/d_P}z,t^{1/d_P}w),
 \qquad d_P=1+a\alpha,                              \tag{22}
\]
scales \(h,P,f_0\) by \(t\) and preserves every cut.
Both the cone total and its inner total are contractible.
Their Mayer--Vietoris sequence says that
\(\widehat J_0\to\widehat M_{{\rm ext},0}\) is a homology
equivalence. Exterior continuation transports this fact to \(s>0\).
Attaching the acyclic \(\widehat W\) shows that
\(\widehat{\mathcal T}_s\) is acyclic.

Now fix \(s>0\) and choose an ordinary Milnor ball \(B_{\rm m}\)
strictly inside all the outer cuts, with value radius
\(0<\delta'\le\delta\). Let \(\widehat N\) be its closed
Milnor disk total, and let \(\widehat A\) be the intervening
part of the large total, with intersection
\(\widehat S=\partial B_{\rm m}\cap f_s^{-1}(D_{\delta'})\).
On \(\widehat A\) and \(\widehat S\), \(f_s\) is a proper
submersion on every face over the whole disk: this uses the
outer boundary proof, the ordinary Milnor-sphere choice,
and uniform isolation in the intervening interior.
Lifting contractions of the disk by face-tangent vector fields
shows that page inclusions
\(A\to\widehat A\) and \(S\to\widehat S\) are homotopy
equivalences. Relative Mayer--Vietoris therefore gives an
integral isomorphism
\[
 H_4(\widehat N,N;\Z)\ \xrightarrow{\;\cong\;}\
 H_4(\widehat{\mathcal T}'_s,\mathcal T'_s;\Z).       \tag{23}
\]
The ordinary total is contractible by the cited Milnor theorem.
Shrinking the large positive value radius from \(\delta\) to
\(\delta'\) is an isotopy: lift the radial field on the
annulus of nonzero values, tangent to all outer faces and
zero near value zero. Thus its disk total stays acyclic.
The relative connecting maps now identify (23) with an
integral page isomorphism \(H_3(N)\cong H_3(\mathcal T_s)\).
An ordinary filling is unchanged inside the larger total.
Transporting its two boundaries back to radius \(\delta\)
along disjoint angular rays adds collars with disjoint
values and no new intersections. Thus (23) also preserves
the ordered filling form.

\section{Exterior coordinate and its signed boundary}

At \(s=0\) take \(c>0\), write \(q=h\), and cut the exterior
at \(\operatorname{Re}q=c/2\). Call its pieces \(A,B\).
On \(A\), \(P=c-q\) stays away from zero; the inner block
is absent. The compact \(P\)-pages trivialize over the
convex allowed value region by the boundary estimates.
Contracting the weighted base and transporting the \(P\)-page
retracts \(A\) to \(Q\).
On \(B\), trivialize \(h\) over its convex regular-value
region, respecting its boundary and the labels in (19).
For
\[
 r(v)=\max((|y|/M)^{1/\alpha},(|z|/M)^{1/\beta},|w|)
\]
scale \(v\) by \(t\) from \(1\) to \(\lambda/r(v)\).
Retain the base-page label and put
\[
 q_t=c-t^{d_P}P(v)=(1-t^{d_P})c+t^{d_P}q.            \tag{24}
\]
This stays in the same convex base region, respects the exterior
cut, and retracts \(B\) to the actual \(J\), relative to \(J\).
The seam is \(\calH\times Q\) times a contractible interval.
Its map to \(J=\calH\times L\) is the identity on the first
factor and weighted radial projection on the second.

Consequently the full page has the homotopy pushout, with
the displayed geometric attachment,
\[
                W\ \longleftarrow\ \calH\times Q\
                   \longrightarrow\ Q.           \tag{25}
\]
In particular its integral Mayer--Vietoris sequence is
\[
 0\longrightarrow K\longrightarrow H_3(\mathcal T_s)
 \xrightarrow{\partial}H\otimes E
 \xrightarrow{\overline{1\otimes\tau}}H/P_*D.       \tag{26}
\]
The corresponding exterior sequence has a single additional
\(H_3(J)=\Z\), generated by a lens fiber. It maps to zero in
\(H_3(W)\), since the individual inner fiber has no third
homology in (20). Thus the exterior image in the full page maps
isomorphically to \(H\otimes\ker\tau\).
Here \(H_2(J)=H/4H\); there is no Kunneth torsion involving
the free group \(H\).

We fix signs at the chain level. Write
\(\mathfrak S(C)=[S^1_y]\times C\), circle first, so that
\(\partial\mathfrak S(C)=-\mathfrak S(\partial C)\).
In (21) the \(K\)-identification and the mod-four injection
\(\Psi\) are \(-\partial_{\calD}\mathfrak S^{-1}\).
In degree two use negative circle-first unsuspension.
Then the base-first torus \([a\times\gamma]\) maps to
\(+a\) in \(H/P_*D\).

For a full cycle write its chains as \(w+e_A+e_B\), with
\(\partial w=b_J\), and define \(t_C=-\partial e_A\).
Then \(\partial e_B=t_C-b_J\).
The prism of the stopped retraction (24) has boundary
\(r_{1*}t_C-t_C\). Adding it to \(e_B\) proves
\[
 z=[t_C],\qquad [b_J]=(1\otimes\tau)z\in H/4H.       \tag{27}
\]
Thus the last arrow of (26) has the positive sign in (27);
it is not left to a convention chosen later.

\section{The geometric root-disk module and the actual transfer}

\begin{lemma}[Braid-natural Bockstein transfer]\label{lem:transfer}
With \(\beta[\overline C]=[\partial C/4]\), for any
\(b\in H/4H=H_2(J)\) and any \(\eta\) with \(\beta\eta=b\),
\[
                       \Psi i_*\eta=T_*b\bmod4D.  \tag{28}
\]
This identity is for the actual graph family and its actual
covering monodromy, not for an arbitrarily substituted cover.
\end{lemma}
\begin{proof}
Work first over a regular \(p\in\calH_c\), with four distinct
roots \(r_1,\ldots,r_4\), all inside \(|w|<\lambda/4\).
The inner fiber \(F_p=\{yz=R_s(p,w)\}\cap Y_\lambda\)
has the collapse (20) with base a disk and root set
\(\{r_1,\ldots,r_4\}\). The anchor \(\gamma_p\) is the
circle (18) at \(\theta=0\). Balancing carries it
orientation-preservingly to the core circle at \(w_0\).
It induces a homotopy equivalence of pairs.
Adjoin the constant-angle disk section to the relative
subspace; this changes no relative homology, since that disk
meets the anchor at one point.
Relative suspension and the endpoint exact sequence now give
\[
 H_2(F_p,\gamma_p)=
 H_1(\Delta,\{w_0,r_1,\ldots,r_4\})
 \cong\{(a_1,\ldots,a_4,a_0):\ \sum a_j+a_0=0\}.
                                                               \tag{29}
\]
Let \(T_i\) be the disk swept over an arc from \(w_0\) to \(r_i\),
with circle-first orientation. Its boundary is \(+\gamma_p\);
its root coordinate is \(+e_i\). Thus (29) is the actual
relative root-disk module \(\Z^4\), its boundary is
\((a_i)\mapsto(\sum a_i)\gamma_p\), and the absolute
spheres are \(\ker\sum\).

The boundary \(L_p=\partial F_p\) is the lens link obtained
by continuing the nonvertex cone boundary. The quotient cover
\((U,V)\mapsto(U^4,V^4,UV)\) identifies one positive \(y\)-turn
with the deck generator \((U,V)\mapsto(iU,-iV)\).
Hence \(H_1(L_p)=\Z/4\), generated by \(\gamma_p\), and
\(H_2(L_p)=0\). The latter also follows from orientability
and the vanishing of \(H^1(L_p;\Z)\).
The exact sequence of \((L_p,\gamma_p)\) supplies a unique
relative class \(S_p\) with \(\partial S_p=4\gamma_p\).

We determine its image in (29), including the possible pure-braid
ambiguity. Move any four distinct roots inside the inner root
disk along a braid and put \(Q_t(w)=\prod_i(w-r_i(t))\).
Choose a disk velocity \(v_t\), supported in that disk, equal
to \(\dot r_i(t)\) in a neighborhood of each root. The function
\[
                 \ell_t=(\dot Q_t+Q_t'v_t)/Q_t
\]
extends smoothly at every root. With a real cutoff \(\chi\)
zero near the \(y\)-face and one near the disjoint \(z\)-face,
the lift
\[
 \dot w=v_t,\qquad
 \dot y=\chi\operatorname{Re}(\ell_t)y,\qquad
 \dot z=(\ell_t-\chi\operatorname{Re}(\ell_t))z       \tag{30}
\]
preserves \(yz=Q_t(w)\), all active boundary faces, and the
\(y\)-angle. It fixes \(w_0,y\) on the anchor, transporting
its \(z\)-coordinate by \(Q_t(w_0)/y\).
The compact flow realizes transport of the full triples
\((F_t,L_t,\gamma_t)\). Under (29) its return is precisely
the permutation of the four endpoint coordinates.
In particular pure braids add no hidden sphere class.

The unique \(S_p\) is invariant under every such return.
Every root permutation can be realized by this auxiliary
deformation of the same fiber triple, so its root coefficients
are equal. Their sum is four. They are consequently all one:
\[
                       S_p\longmapsto
                  N=(1,1,1,1).                    \tag{31}
\]
This use of auxiliary braids computes a single fiber class;
it does not assert that the actual graph has full symmetric
monodromy.

Now take an oriented loop \(a\) in the actual \(\calH_c\)
avoiding the branch values. The actual polynomials
\(R_s(p,\cdot)\) induce a particular braid, to which (30)
applies. Sweep a representative \(S_0\) around the cut loop,
obtaining an interval-first chain \(T_a\) with endpoint \(S_1\).
The anchor is fixed at the seam. Since \(H_2(L)=0\), choose
a seam cap \(C_a\) with \(\partial C_a=S_1-S_0\).
The product boundary formula gives
\[
 \partial T_a=S_1-S_0-4(a\times\gamma),\qquad
 \eta_a=-[(T_a-C_a)\bmod4],\qquad
 \beta\eta_a=[a\times\gamma].                       \tag{32}
\]
Changing the cap adds the reduction of an integral lens
three-cycle, which dies in \(W\). The same is true of the
ambiguity in an arbitrary Bockstein lift on \(J\).

Slice away from the seam. For an oriented small base interval \(I_a\)
and an arc \(I_i\) from \(w_0\) to \(r_i\),
\[
 -I_a\times T_i=\mathfrak S(I_a\times I_i),\qquad
 -\partial_{\calD}(I_a\times I_i)=+[I_{a,i}].
                                                               \tag{33}
\]
Thus (31)--(33) give the positive diagonal vector for
\(\Psi i_*\eta_a\), with exactly the sign in (28).
For the actual root permutation \(\sigma\), the components of
\(\pi^{-1}(a)\) correspond to its orbits. Slicing their oriented
fundamental cycles gives the orbit-sum vectors in \((\Z/4)^4\).
These have disjoint supports, so the slice map on
\(H_1(\pi^{-1}(a);\Z/4)\) is injective, including orbits of
length two or four. The full actual transfer has the same
diagonal vector (31). Equality therefore holds already over
that actual covering loop, not merely in a rational module.
Naturality under inclusion into the full root page proves (28).
Such loops generate \(H\).
\end{proof}

The compact integral matrix
\[
 B=\begin{pmatrix}
 1&0&0&1\\-1&1&0&1\\0&-1&1&1\\0&0&-1&1
 \end{pmatrix}                                    \tag{34}
\]
has columns \(e_1-e_2,e_2-e_3,e_3-e_4,N\) in (29).
Its Smith diagonal is \((1,1,1,4)\).
This certifies the index-four quotient of the actual relative
root-disk group by the absolute spheres and the lens cap.
It is not a substitute for the full page lattice.
The retained Jacobian output gives explicit unimodular left
and right transformations; \texttt{verify.py} replays them
with exact integers and separately checks every one of the
24 possible root permutations over \(\Z/4\).

\section{Exact integral image, not just an index bound}

For \(x\in H_3(\mathcal T_s)\), put \(z=\partial x\).
Since \(4E\subset\ker\tau\), the exterior image supplies a
unique class \(x_{\rm ext}(4z)\) with exterior coordinate \(4z\).
Define
\[
             u=4x-x_{\rm ext}(4z)\in K,\qquad
             x\longmapsto(u/4,z).                  \tag{35}
\]
The map is injective: if \(z=0\), (26) puts \(x\) in \(K\)
and (35) is \((x,0)\).
Both \(K\) and the image of \(\partial\) are free, so (26)
also excludes hidden torsion in \(H_3(\mathcal T_s)\).

Write \(x=w+v\), with \(\partial w=b_J\),
\(\partial v=-b_J\). By (27), \([b_J]=b=(1\otimes\tau)z\).
Choose an integral chain \(A\) in \(J\) with
\(\partial A=4b_J\). The two closed chains are
\[
                  w_4=4w-i_WA,\qquad
                  v_4=4v+i_{\rm ext}A.             \tag{36}
\]
Their sum is \(4x\), and \(v_4\) represents
\(x_{\rm ext}(4z)\) in the full page. Reducing the first
chain modulo four, applying (28), and using
\(\beta[A\bmod4]=b\), gives the necessary residue
\[
                        u\equiv-T_*b\pmod{4D}.     \tag{37}
\]
The minus sign comes from subtracting the interface cap,
not from reversing a chosen transfer.

For any integral lift \(b_H\in H\), set
\(l=(u+T_*b_H)/4\). Equation (37) says \(l\in D\);
then (3) gives \(P_*l=b_H\) and \(\pr(l)=u/4\).
This proves containment in (5).

Conversely, let \(l,z\) satisfy (5). Since \(4H\subset P_*D\),
the last arrow of (26) kills \(z\). Choose an actual page class
\(x\) over \(z\), with numerator \(u_x\). Put
\(u_l=4l-T_*P_*l\in K\).
Equation (37) gives
\[
                  u_x-u_l\in K\cap4D=4K.           \tag{38}
\]
The equality is integral: if \(4d\in K\), then \(4P_*d=0\)
in the torsion-free \(H\), so \(d\in K\).
Adding \((u_l-u_x)/4\in K\) to \(x\) realizes exactly
\((\pr(l),z)\). This proves surjectivity onto all of (5).
It uses the actual geometric exact sequence and the actual
transfer, not equal ranks or an abstract isometric lattice.

\section{Both orders of the filling form}

\subsection{Cross terms and the exterior summand}

Inner cycles bound in the acyclic \(\widehat W\).
An exterior page representative of \(H\otimes\ker\tau\)
need not initially bound in \(\widehat M_{\rm ext}\).
However, the interface disk total continues to the product
of the acyclic base total with \(L\); hence
\[
 H_3(\widehat J)=H_3(\widehat M_{\rm ext})=\Z .
\]
The isomorphism is the actual inclusion, as proved using (22).
The page lens-fiber generator maps to a generator of this
group and dies in the full page. Subtracting the requisite
integer multiple of it from an exterior representative
makes that representative bound in \(\widehat M_{\rm ext}\).
This adjustment can be transported over a short value arc.
Put the inner and exterior fillings on their respective sides
of the interface bicollar. They are disjoint. This works
with the inner filling first and also with the exterior
filling first, proving both ordered cross terms vanish.
Since \(4E\subset\ker\tau\), these exterior classes span
the whole rational exterior summand.

Continue the exterior fillings to \(s=0\) and lift them
fully under the cone cover. All these fillings stay off
the vertex. Their intersection number upstairs is four
times the downstairs number, since the cover is
orientation-preserving and unramified there.
Embed the lifted fillings in the smooth weighted tube
of \(h+\widetilde P\). Its disk total is contractible
by the simultaneous positive weights, so replacing
fillings with the same boundary changes no intersection
number.

Here is the product computation without invoking an external
Seifert-product theorem. For loops \(a,v\) in the two plane
pages, the join of their weighted radial transports at
values \(tc,(1-t)c\) is the class whose seam coordinate
is \(a\otimes v\). The product of the two radial cone
fillings, restricted to the triangle \(t,t'\ge0,\ t+t'\le1\),
bounds this join; the edges at a collapsed loop contribute
no chain boundary. Construct the second filling on the
positive neighboring ray. Their intersections are localized
where both factor values vanish and, after relative
general-position perturbation, their local numbers are
products of the two plane filling intersection numbers.
Moving the second real two-dimensional factor past the
first real two-dimensional factor contributes
\((-1)^{2\cdot2}=+1\). A common reversal of the join
orientation occurs in both fillings and cancels.
Thus the upstairs ordered form is
\(s_h\otimes\widetilde s\).
The seam coordinate of a full lift is
\((1\otimes\Tr)z\). Dividing intersections by four gives
exactly \(s_h\otimes e\) on the exterior rational summand.

\subsection{The inner summand and its negative normal determinant}

Work in the inner \(h\)-model over the whole value disk.
Let \(\widehat G\) be the acyclic base total times the
\(w\)-disk and \(\widehat{\calD}\) its smooth root total.
For \(l_1,l_2\in K\) on the reference and positive neighboring
root pages, choose surface fillings
\(\Sigma_i\subset\widehat{\calD}\) with
\(\partial\Sigma_i=l_i\).
The kernel condition gives surface chains
\(V_i\subset G_{c_i}\) with \(\partial V_i=-l_i\).
Since \(H_2(\widehat G)=0\), choose three-chains \(\Gamma_i\)
with
\[
                      \partial\Gamma_i=\Sigma_i+V_i.             \tag{39}
\]
Prescribe a collar of \(\Sigma_i\) with positive real
normal coordinate \(R_s=t\). Near \(l_i\), the normal
can be chosen tangent to \(h=c_i\), since the root page
is regular. Extend the remaining chains relatively.
Make the \(\Sigma_i\) transverse in the real four-dimensional
root total. Generic perturbations of the unprescribed
parts of \(\Gamma_i\) meet the root in one-dimensional
chains \(I_i\), disjoint from the other \(\Sigma_j\)
and from each other: their dimension sums in the root
total are \(1+2<4\) and \(1+1<4\).
The distinct page-value collars separate all boundary contacts.

Choose distinct positive balancing parameters \(\sigma_1,\sigma_2\)
close enough to the value used in (20). Sweep \(\Gamma_i\)
by the \(y\)-circle at
\[
 |y|=\sigma_i\sqrt{|R_s|},\quad
 |z|=\sigma_i^{-1}\sqrt{|R_s|},\quad z=R_s/y .
\]
Give the sweep orientation \(-\mathfrak S(\Gamma_i)\).
Then
\[
 \partial(-\mathfrak S(\Gamma_i))=\mathfrak S(V_i),
 \qquad -\partial_{\calD}V_i=l_i,                  \tag{40}
\]
so these are fillings of the positive \(K\)-classes in our
fixed convention. They are legitimate integral mapped
chains, not merely pictures of smooth submanifolds.
Near a transverse incidental root curve the collapsed
domain is \(I_i\times\operatorname{cone}(T^2)\).
The closed oriented torus link makes the top-chain
boundary cancel; the cone tip is not a codimension-one
face. At a prescribed root collar the collapse is an
ordinary disk. These observations justify (40) after a
compatible triangulation.

Off the root the two swept chains are disjoint, because
their balancing radii differ. Their only common root
contacts are \(\Sigma_1\cap\Sigma_2\).
The incidental cone tips are disjoint from the other
chain, so local general-position approximations there
change neither the class nor any intersection number.
At a prescribed contact, the normal disks in complex
coordinates \(y,z\) are
\[
                         z=\sigma_i^{-2}\overline y.
\]
Put \(c_i=\sigma_i^{-2}\). In real coordinates their
ordered intersection determinant is
\[
 \det\begin{pmatrix}
 1&0&1&0\\0&1&0&1\\c_1&0&c_2&0\\0&-c_1&0&-c_2
 \end{pmatrix}
                    =-(c_2-c_1)^2<0.              \tag{41}
\]
The common sweep reversal cancels between the two disks.
Exchanging a real two-dimensional disk with a real
two-dimensional root filling contributes \(+1\).
Thus every root filling intersection is reversed, in
the fixed neighboring-page order:
\[
                  s_{\rm inner}(l_1,l_2)=-s_d(l_1,l_2).            \tag{42}
\]
The same determinant is negative in the reverse order;
no symmetry of the plane Seifert form was assumed.
Equations (23), (35)--(38), and (42), together with the
exterior computation, prove Theorem~\ref{thm:main}.

\section{Exact replay, interpretation, and remaining frontier}

Run \texttt{python3 -B research/m45\_seifert/verify.py} from
the campaign root. It uses exact arithmetic, not numerical
root tracking, and replays the following compact identities:
\begin{itemize}
\item the source polynomial and scaling identities, the gradient
Bezout identity (8), and the polar polynomial (10);
\item the geometrically realized root-disk matrix (34), including
the retained Jacobian unimodular transformations and index four;
\item orbit-slice injectivity, the norm vector, and full-transfer
degree for all 24 four-sheet monodromy permutations over \(\Z/4\);
\item the chain boundaries (32), (36), their mod-four minus sign,
and the integral saturation argument underlying (38);
\item the normal determinant (41), its reversed order, and the
positive exterior product shuffle sign.
\end{itemize}
The accompanying certificate records exactly which assertions are
computations and which require the geometric proof. It makes no
claim that enumerating permutations proves the existence of the
local comparison, or that a successful Smith calculation identifies
the actual Milnor lattice on its own. Actual realization is supplied
by (18)--(20), (23)--(25), and (29)--(33).

As a consistency check, with \(m=4ab+a+b\) the cone page has
rank \(m+1\). Indeed the isolated weighted plane polynomial
\(\widetilde P\) has weights \((b,a)\), degree \(m\), and
Jacobian length
\((m-b)(m-a)/(ab)=(4a+1)(4b+1)\); the length follows
from the two-generator weighted Koszul Hilbert series.
Its free degree-four page quotient has Euler characteristic
\((1-\mu(\widetilde P))/4=-m\).
Equation (5) therefore has rank
\((51-6)+6(m+1)=24ab+6a+6b+51\), agreeing with
\texttt{C059APPLICATION}. This equality is a check after the
integral proof, not a surjectivity argument.

The mathematical improvement over the strongest prior result is
the integral local comparison, the full realized mod-four
extension, and the ordered form, together with a uniform
outer-isolation proof. No isometry between the two resulting
\(\Lambda_{\pi_i}\) is asserted here. M45's arithmetic congruence
and classification applications, all independent M23 obligations,
independent review of this proof, and any claim of publishable
novelty remain open. No Reviewer verdict is self-recorded.

\begin{thebibliography}{9}
\bibitem{Massey}
D. B. Massey, \emph{Non-isolated Hypersurface Singularities and
L\^e Cycles}, arXiv:1410.3312v2, November 13, 2015.
Theorem 1.8, Remark 1.9, and Section 1, p.~6, items (2)--(5).
\end{thebibliography}
\end{document}
